\chapter{Java}
\label{cap:java-teoria}

\begin{edital}
Java (versão 6+), JavaEE (6+), JakartaEE, JPA (2+), frameworks JUnit,
Hibernate, JSF, Primefaces, Spring/SpringCloud/SpringBoot. É a
linguagem-base do perfil.
\end{edital}

\section{Exceções: \textit{checked} x \textit{unchecked}}

\begin{conceito}
\begin{center}
\begin{tabular}{@{}p{2.6cm}p{4.6cm}p{4.6cm}@{}}
\toprule
& \textbf{Checked (verificada)} & \textbf{Unchecked (não verificada)} \\
\midrule
Superclasse & \texttt{Exception} (não \texttt{RuntimeException}) &
\texttt{RuntimeException} \\
Compilador & \textbf{exige} \texttt{try-catch} \textbf{ou}
\texttt{throws} & não exige nada \\
Exemplos & \texttt{IOException}, \texttt{SQLException} &
\texttt{NullPointerException}, \texttt{ArrayIndexOutOfBounds},
\texttt{IllegalArgument} \\
\bottomrule
\end{tabular}
\end{center}
A ideia por trás da divisão: exceções \textit{checked} representam
condições que um programa \textbf{razoável deveria antecipar e tratar}
(um arquivo pode não existir, uma conexão de rede pode cair) --- o
compilador força a lidar com elas. Exceções \textit{unchecked} tipicamente
representam \textbf{erros de programação} (acessar índice inválido,
referência nula) que, em teoria, um código correto nunca deveria disparar
--- por isso não são forçadas no nível do compilador.
\end{conceito}

\begin{cuidado}
\texttt{NullPointerException} é \textbf{unchecked} (estende
\texttt{RuntimeException}); \texttt{IOException} é \textbf{checked}. Uma
exceção \textit{unchecked} \textbf{não} exige \texttt{throws} obrigatório
--- é o contrário do que soa intuitivo pelo nome.
\end{cuidado}

\section{Polimorfismo: \textit{overload} x \textit{override}}

\begin{conceito}
\begin{center}
\begin{tabular}{@{}p{2.6cm}p{4.6cm}p{4.6cm}@{}}
\toprule
& \textbf{Overload (sobrecarga)} & \textbf{Override (sobrescrita)} \\
\midrule
Onde & mesma classe & subclasse \\
Assinatura & \textbf{muda} (parâmetros diferentes) & \textbf{igual} à do
pai \\
Resolução & \textit{compile-time} (estática) & \textit{runtime}
(dinâmica) \\
Anotação & --- & \texttt{@Override} \\
\bottomrule
\end{tabular}
\end{center}

A \textbf{assinatura} de um método é o nome mais a lista de parâmetros ---
o \textbf{tipo de retorno não entra nela}. Consequência direta: dois
métodos no mesmo escopo que diferem \textbf{apenas} no tipo de retorno
(\texttt{int calc(int x)} e \texttt{long calc(int x)}) \textbf{não
compilam} --- o compilador não os reconhece como sobrecarga válida, é
declaração duplicada da mesma assinatura.
\end{conceito}

\begin{cuidado}
Sobrecarga = mesmo nome, assinatura diferente, decidida em compilação.
Sobrescrita = mesma assinatura, na subclasse, decidida em execução.
Sobrescrever \textbf{não} viola Liskov por si só (ver
Seção~\ref{sec:liskov-teoria}) --- viola se quebrar o contrato esperado
pelo cliente, não pelo simples fato de existir.
\end{cuidado}

\section{Orientação a objetos e SOLID}

\begin{conceito}[Os quatro pilares: o mesmo propósito, alvos diferentes]
Os quatro pilares clássicos (abstração, encapsulamento, herança,
polimorfismo) têm o \textbf{mesmo} propósito: limitar o que o resto do
sistema precisa saber sobre uma parte dele. O que muda entre eles é
\textbf{o que} exatamente está sendo escondido:

\begin{center}
\begin{tabular}{@{}p{2.7cm}p{4.1cm}p{4.5cm}@{}}
\toprule
\textbf{Pilar} & \textbf{Esconde\ldots} & \textbf{Como aparece no código} \\
\midrule
Abstração & os detalhes irrelevantes para quem usa & a classe expõe
\textit{o que} faz (\texttt{sacar}), não \textit{como} \\
Encapsulamento & o \textbf{estado} interno & atributo \texttt{private} +
operações públicas que \textbf{protegem a regra} de negócio \\
Herança & a parte comum, escrita uma única vez & \texttt{extends}:
relação \textbf{é-um} \\
Polimorfismo & \textbf{qual} implementação vai atender a chamada &
\texttt{conta.sacar()} executa o método da classe \textbf{real} do
objeto, decidido em \textbf{tempo de execução} \\
\bottomrule
\end{tabular}
\end{center}

Duas confusões comuns sobre esses pilares. \textbf{Encapsulamento não é
"gerar getter e setter para todo atributo"}: um \texttt{setSaldo} público
expõe o estado tão livremente quanto um atributo público faria, e a regra
de negócio ("não sacar além do saldo") fica sem dono. Encapsular é expor
\textbf{comportamento} controlado, não simplesmente esconder o campo atrás
de dois métodos que fazem a mesma coisa que acesso direto faria. O
\textbf{polimorfismo que importa} é o de \textbf{sobrescrita}, resolvido
em \textit{runtime} pela classe real do objeto --- a sobrecarga é
resolvida em \textbf{compilação}, pelos tipos declarados, e por isso
alguns autores nem a contam como polimorfismo "de verdade".

Herança é o pilar com \textbf{contraindicação}: ela acopla a subclasse à
implementação interna do pai. Quando a relação real é "\textbf{tem-um}" e
não "é-um", a escolha correta é \textbf{composição} (Capítulo
\ref{cap:padroes-uml-teoria}).
\end{conceito}

\begin{conceito}[SOLID --- os cinco princípios]
\begin{center}
\begin{tabular}{@{}p{1.2cm}p{9.6cm}@{}}
\toprule
\textbf{Letra} & \textbf{Princípio} \\
\midrule
S & \textit{Single Responsibility} --- uma responsabilidade por classe \\
O & \textit{Open/Closed} --- aberta a extensão, fechada a modificação \\
L & \textbf{Liskov} --- subtipo substitui o tipo base sem quebrar o
programa \\
I & \textit{Interface Segregation} --- interfaces pequenas e específicas \\
D & \textit{Dependency Inversion} --- depender de abstração, não de
implementação concreta \\
\bottomrule
\end{tabular}
\end{center}
\end{conceito}

\subsection{Liskov em detalhe: a violação mais cobrada}
\label{sec:liskov-teoria}

\begin{exemplo}[Retângulo e Quadrado --- a herança que compila e ainda assim quebra]
\begin{center}
\begin{tabular}{@{}p{5.5cm}p{6.1cm}@{}}
\toprule
\textbf{O código} & \textbf{O que acontece} \\
\midrule
\texttt{class Retangulo \{ setLargura; setAltura; area() \}} & quem usa
esta classe assume: mudar a largura não muda a altura \\
\texttt{class Quadrado extends Retangulo} --- e cada \textit{setter} muda
\textbf{os dois} lados, para manter o quadrado quadrado & a classe está
coerente consigo mesma\ldots \\
\texttt{r.setLargura(5); r.setAltura(4);} \newline
\texttt{assert r.area() == 20;} & \ldots{} mas com um \texttt{Quadrado} no
lugar do \texttt{Retangulo}, a área dá \textbf{16}. O código do cliente,
que nunca soube da existência do quadrado, quebrou \\
\bottomrule
\end{tabular}
\end{center}

Liskov não fala sobre a subclasse isoladamente: fala sobre \textbf{quem
consome o supertipo}. O critério prático: a subclasse não pode
\textbf{exigir mais} do que o pai exigia (fortalecer uma pré-condição),
não pode \textbf{entregar menos} do que o pai entregava (enfraquecer uma
pós-condição), e não pode quebrar uma \textbf{invariante} que o cliente já
podia supor como verdadeira. O exemplo canônico do outro lado é o
\texttt{Pinguim} que herda de \texttt{Ave} com um método \texttt{voar()}:
a hierarquia parece natural biologicamente, mas o contrato de código não
fecha.
\end{exemplo}

\begin{cuidado}
O próprio compilador do Java já barra \textbf{três} violações de Liskov
na sobrescrita: o método sobrescrito não pode \textbf{reduzir a
visibilidade} em relação ao pai, não pode declarar uma exceção
\textbf{checked mais ampla} que a do pai, e o tipo de retorno precisa ser
o mesmo ou um \textbf{subtipo} dele (covariante). O que o compilador
\textbf{não} consegue barrar é a quebra de contrato \textbf{semântico}
(como no exemplo do Retângulo/Quadrado) --- e é exatamente ali que a
prova cobra o entendimento real do princípio.
\end{cuidado}

\begin{cuidado}
ISP (interface enxuta) e SRP (uma responsabilidade por classe) não são o
mesmo princípio, apesar de soarem parecidos --- ISP fala de
\textbf{interface}, SRP fala de \textbf{classe}.
\end{cuidado}

\section{Interface x classe abstrata}

\begin{conceito}
\begin{center}
\begin{tabular}{@{}p{3.2cm}p{4.2cm}p{4.2cm}@{}}
\toprule
& \textbf{Classe abstrata} & \textbf{Interface} \\
\midrule
Construtor & \textbf{tem} & \textbf{não tem} \\
Atributo de instância & tem (com estado) & só \texttt{public static
final} (constante) \\
Herança múltipla & \textbf{não} --- só uma classe pode ser estendida &
\textbf{sim} --- várias podem ser implementadas \\
Método concreto & sim, desde sempre & sim, a partir do Java 8
(\texttt{default} e \texttt{static}) \\
Palavra-chave & \texttt{extends} & \texttt{implements} \\
\bottomrule
\end{tabular}
\end{center}

Regra de bolso: uma classe \textbf{estende uma} classe (abstrata ou não) e
\textbf{implementa várias} interfaces. Use classe abstrata quando há
\textbf{estado e comportamento comuns} genuínos a compartilhar entre
subclasses; use interface quando o que importa é apenas o
\textbf{contrato} a ser cumprido.
\end{conceito}

\begin{cuidado}
Desde o Java 8, interfaces têm métodos \texttt{default} com corpo
implementado --- "interface não pode ter implementação" está
desatualizado e é falso hoje. O que a interface continua \textbf{sem
ter} é \textbf{construtor} e \textbf{atributo de instância} --- é por aí
que a distinção real se sustenta. Java não tem herança múltipla de
\textbf{classes} (só de \textbf{tipo}, via múltiplas interfaces
implementadas).
\end{cuidado}

\section{Coleções (\textit{Collections})}

\begin{conceito}[Contrato x implementação --- por que o framework é desenhado assim]
\texttt{List}, \texttt{Set} e \texttt{Map} são \textbf{interfaces} ---
dizem que \textbf{garantia} você tem (ordem, unicidade, associação
chave-valor) e nada sobre \textbf{como} ela é entregue internamente.
\texttt{ArrayList} e \texttt{LinkedList} são duas \textbf{implementações}
diferentes da mesma interface \texttt{List}, com os custos opostos das
estruturas de dados vistas no Capítulo \ref{cap:programacao-teoria}
(Seção~\ref{sec:estruturas-dados-teoria}). Por isso a boa prática é
declarar pela interface (\texttt{List<String> lista = new
ArrayList<>()}): trocar a implementação depois não exige mexer em quem a
usa.

\begin{center}
\begin{tabular}{@{}p{2.2cm}p{3.8cm}p{4cm}@{}}
\toprule
\textbf{Interface} & \textbf{Implementações} & \textbf{Característica} \\
\midrule
\texttt{List} & \texttt{ArrayList}, \texttt{LinkedList} & ordenada,
permite duplicatas \\
\texttt{Set} & \texttt{HashSet}, \texttt{TreeSet}, \texttt{LinkedHashSet}
& sem duplicatas \\
\texttt{Map} & \texttt{HashMap}, \texttt{TreeMap}, \texttt{LinkedHashMap}
& associação chave-valor \\
\bottomrule
\end{tabular}
\end{center}

\texttt{ArrayList} (array dinâmico, acesso $O(1)$ por índice) x
\texttt{LinkedList} (lista ligada, inserção/remoção $O(1)$ no meio se já
se tem o nó). \texttt{HashMap} (sem ordem garantida) x
\texttt{LinkedHashMap} (preserva ordem de inserção) x \texttt{TreeMap}
(mantém ordenação pela chave). \texttt{getOrDefault} retorna o valor
\textbf{padrão} fornecido se a chave não existir --- não lança exceção
nem retorna um código de erro. \texttt{==} compara \textbf{referência}
(mesmo objeto na memória); \texttt{.equals()} compara \textbf{conteúdo}.
\end{conceito}

\begin{conceito}[O contrato equals/hashCode: por que um objeto "desaparece" de um HashSet]
Coleções com "Hash" no nome não procuram um elemento comparando com
\textbf{todos} os demais --- elas calculam o \texttt{hashCode()} do
objeto para achar o \textbf{compartimento} correto, e só dentro dele usam
\texttt{equals()} para confirmar a igualdade exata. Isso impõe um
contrato obrigatório:

\begin{center}
se \texttt{a.equals(b)} é \texttt{true}, então \texttt{a.hashCode() ==
b.hashCode()} \textbf{obrigatoriamente}
\end{center}

A recíproca \textbf{não} vale --- dois objetos diferentes podem
compartilhar o mesmo \textit{hash} (é colisão, e é um evento normal e
esperado).

A consequência prática cobrada em prova: sobrescrever \texttt{equals} e
\textbf{esquecer} de sobrescrever \texttt{hashCode} junto produz um bug
silencioso. Dois objetos considerados "iguais" pelo \texttt{equals} caem
em compartimentos \textbf{diferentes} do hash (porque o \texttt{hashCode}
herdado de \texttt{Object} usa identidade de memória), então o
\texttt{HashSet} aceita a duplicata sem reclamar, e
\texttt{contains}/\texttt{get} de um \texttt{HashMap} \textbf{não
encontra} um objeto que está fisicamente lá dentro. Nada estoura --- só dá
a resposta errada silenciosamente, o pior desenho possível para um item de
"qual a saída deste código".

\texttt{TreeSet} e \texttt{TreeMap} \textbf{não} usam \texttt{hashCode}
--- eles ordenam os elementos, e por isso dependem de
\texttt{compareTo()} (via \texttt{Comparable}) ou de um
\texttt{Comparator} explícito. Uma classe sem nenhum dos dois dispara
\texttt{ClassCastException} ao ser inserida num \texttt{TreeSet} --- e ao
mesmo tempo funciona perfeitamente bem num \texttt{HashSet}.
\end{conceito}

\subsection{String: imutabilidade e StringBuilder}

\begin{conceito}
\texttt{String} é \textbf{imutável}: toda operação que parece "alterar"
(\texttt{concat}, \texttt{+}, \texttt{replace}, \texttt{toUpperCase})
na verdade devolve um \textbf{objeto novo}, deixando o original intacto.
Por isso, concatenar dentro de um laço cria um objeto novo descartado a
cada iteração --- em concatenação intensiva, usa-se
\texttt{StringBuilder} (mutável, não sincronizado, mais rápido) ou
\texttt{StringBuffer} (mutável e sincronizado entre threads, mais lento).

Literais de string iguais compartilham o \textbf{\textit{string pool}}
(uma área especial de memória), então \texttt{"a" == "a"} costuma dar
\texttt{true}; já \texttt{new String("a") == "a"} dá \texttt{false},
porque o \texttt{new} força a criação de um objeto \textbf{fora} do pool,
mesmo com o mesmo conteúdo. Comparação de conteúdo deve sempre usar
\texttt{.equals()}, nunca \texttt{==}.
\end{conceito}

\begin{cuidado}
"\texttt{String} é mutável porque posso reatribuir a variável" é falso:
o que muda ao reatribuir é a \textbf{referência} guardada na variável, não
o objeto \texttt{String} original, que permanece intacto na memória.
\texttt{StringBuilder} (não sincronizado, mais rápido) e
\texttt{StringBuffer} (sincronizado, mais lento) não são intercambiáveis.
\end{cuidado}

\section{JVM, JPA e Java EE}

\begin{conceito}[JVM: por que existe, e as duas áreas de memória]
O compilador Java não gera código de máquina diretamente --- gera
\textbf{bytecode} (\texttt{.class}), uma instrução intermediária que
nenhuma CPU física executa nativamente. Quem executa é a \textbf{JVM}, e
existe uma implementação de JVM para cada sistema operacional --- é
exatamente daí que vem o lema \textit{write once, run anywhere}: o
\texttt{.class} é idêntico em Windows e Linux, o que muda é o intérprete
que o executa.

A separação de memória explica boa parte das perguntas sobre o tema:

\begin{center}
\begin{tabular}{@{}p{2.2cm}p{9.4cm}@{}}
\toprule
\textbf{Área} & \textbf{O que guarda} \\
\midrule
\textbf{Stack} & uma por \textbf{thread}: um quadro por chamada de
método, com parâmetros, variáveis locais e \textbf{referências}. O
quadro morre quando o método retorna \\
\textbf{Heap} & os \textbf{objetos} em si, compartilhados por todas as
threads. Sobrevivem ao método que os criou, enquanto alguma referência
ainda os alcançar \\
\bottomrule
\end{tabular}
\end{center}

A variável local mora na \textit{stack}; o objeto que ela referencia mora
no \textit{heap}. É por isso que "passar o objeto" para um método é, na
prática, passar uma cópia da referência (Capítulo
\ref{cap:programacao-teoria}) --- e é por isso que a \textit{stack} é
literalmente uma pilha (a mesma estrutura LIFO da Seção
\ref{sec:estruturas-dados-teoria}, agora funcionando como mecanismo de
execução): recursão sem fim estoura a pilha
(\texttt{StackOverflowError}); criar objetos demais estoura o heap
(\texttt{OutOfMemoryError}).

\textbf{Coleta de lixo (\textit{garbage collection})}: o coletor libera o
que \textbf{não é mais alcançável} a partir das raízes (variáveis vivas
nas pilhas de todas as threads, campos estáticos) --- não simplesmente o
que "não está mais sendo usado" no sentido intuitivo. Consequências
cobradas: Java \textbf{não tem destrutor determinístico} (não se sabe
\textbf{quando} exatamente um objeto será coletado) e
\texttt{System.gc()} é apenas uma \textbf{sugestão} à JVM, não uma ordem
imediata. Um recurso externo (arquivo, conexão de rede) precisa ser
fechado explicitamente, com \texttt{try-with-resources} --- o coletor de
lixo cuida de memória, não de recursos externos (\textit{handles}).
\end{conceito}

\begin{conceito}[JPA x Hibernate, EAGER x LAZY, problema N+1]
\textbf{JPA} é a \textbf{especificação} de persistência (a interface
padronizada); \textbf{Hibernate} é uma \textbf{implementação} concreta
dessa especificação (o motor de ORM que faz o trabalho de fato).
\texttt{EAGER} (carregamento antecipado, junto com a entidade principal)
x \texttt{LAZY} (carregamento sob demanda, só quando a relação é
efetivamente acessada).

O \textbf{problema N+1}: buscar uma lista de $N$ entidades e, para cada
uma, disparar uma consulta \textbf{separada} para carregar sua relação
associada --- resultando em $N+1$ consultas ao banco em vez de uma só.
Resolve-se combinando \texttt{LAZY} (não carrega a relação por padrão)
com \texttt{JOIN FETCH} explícito na consulta específica que realmente
precisa da relação naquele momento --- trazendo tudo numa única consulta
otimizada, só quando necessário.
\end{conceito}

\begin{cuidado}
Diante do problema N+1, "usar EAGER em tudo" ou "deixar o JPA decidir
sozinho" não são a solução --- a resposta correta é \texttt{LAZY} como
padrão, combinado com \texttt{JOIN FETCH} pontual na consulta que precisa
da relação.
\end{cuidado}

\begin{conceito}[Java EE / Jakarta EE]
\textbf{Java EE/Jakarta EE} reúne especificações como EJB, JPA, JMS e
JSF. \textbf{Jakarta} é a continuação do Java EE sob governança da
Eclipse Foundation (o projeto migrou da Oracle) --- a mudança visível mais
citada é o \textbf{namespace}: \texttt{jakarta.*} substituiu
\texttt{javax.*} nos pacotes.
\end{conceito}

\subsection{Anotações que aparecem em leitura de código}

\begin{conceito}
\begin{center}
\begin{tabular}{@{}p{3.2cm}p{8cm}@{}}
\toprule
\textbf{Anotação} & \textbf{Papel} \\
\midrule
\multicolumn{2}{@{}l}{\textit{Spring --- estereótipos de componente}}\\
\texttt{@Component} & bean genérico gerenciado pelo contêiner \\
\texttt{@Service} & camada de regra de negócio \\
\texttt{@Repository} & camada de \textbf{acesso a dados}; acrescenta
\textbf{tradução automática de exceções} de persistência para a
hierarquia de exceções do Spring \\
\texttt{@Controller} / \texttt{@RestController} & camada web; o
\texttt{@RestController} já embute \texttt{@ResponseBody} \\
\texttt{@Autowired} & injeta a dependência automaticamente \\
\midrule
\multicolumn{2}{@{}l}{\textit{JPA --- mapeamento objeto-relacional}}\\
\texttt{@Entity} & marca a classe como entidade persistente \\
\texttt{@Id} & identifica a chave primária (com \texttt{@GeneratedValue}
para geração automática) \\
\texttt{@Table} / \texttt{@Column} & renomeiam tabela e coluna
(opcionais, senão usam o nome do campo/classe) \\
\texttt{@ManyToOne} & \textbf{muitos} registros desta entidade para
\textbf{um} da outra (é o lado dono da relação, que guarda a chave
estrangeira) \\
\texttt{@OneToMany} & o inverso do anterior; costuma vir acompanhado de
\texttt{mappedBy} \\
\texttt{@OneToOne}, \texttt{@ManyToMany} & relação um-para-um e
muitos-para-muitos \\
\bottomrule
\end{tabular}
\end{center}
O par mínimo para uma entidade JPA funcionar é \texttt{@Entity} +
\texttt{@Id} --- todo o resto das anotações é opcional.
\end{conceito}

\begin{cuidado}
A cardinalidade de \texttt{@ManyToOne} se lê do lado de quem a declara
para o outro lado: \textbf{muitos} pedidos para \textbf{um} cliente ---
inverter essa leitura é o erro mais comum. Atenção também a nomes de
anotação inventados que soam plausíveis mas não existem
(\texttt{@Persistent}, \texttt{@PrimaryKey}) --- o nome exato importa.
\end{cuidado}

\section{Recursos modernos (Java 17/21 LTS)}

\begin{conceito}[Sealed, threads virtuais, streams]
\textbf{\texttt{sealed}/\texttt{non-sealed}/\texttt{final}}: uma classe
\texttt{sealed} exige que suas subclasses sejam explicitamente declaradas
(via \texttt{permits}, ou estando no mesmo arquivo), e cada subtipo deve
ser marcado como \texttt{final}, \texttt{sealed} ou \texttt{non-sealed}
(que reabre a hierarquia para herança livre a partir dali).

\textbf{Threads virtuais (\textit{Project Loom}) x threads de
plataforma}: milhares de threads \textbf{virtuais} compartilham um número
muito menor de \textbf{\textit{carrier threads}} de plataforma (não é uma
relação um-para-um) --- brilham especialmente em cenários dominados por
espera de I/O.
\end{conceito}

\begin{conceito}[Stream: um pipeline preguiçoso]
Um \textit{stream} \textbf{não é uma coleção} --- é um \textbf{fluxo de
processamento} sobre uma fonte de dados, com três partes:

\begin{center}
fonte (\texttt{lista.stream()}) $\to$ operações \textbf{intermediárias}
$\to$ operação \textbf{terminal}
\end{center}

Operações \textbf{intermediárias} (\texttt{filter}, \texttt{map},
\texttt{sorted}, \texttt{distinct}, \texttt{limit}) devolvem \textbf{outro
stream} e são \textbf{preguiçosas} --- nenhuma delas executa no momento
em que é escrita no código. Operações \textbf{terminais}
(\texttt{collect}, \texttt{forEach}, \texttt{reduce}, \texttt{count},
\texttt{anyMatch}) \textbf{disparam} a execução de todo o pipeline
acumulado e o encerram.

Três consequências diretas dessa lógica preguiçosa: um pipeline
\textbf{sem operação terminal não executa nada} (nem o \texttt{map} chega
a rodar); o stream \textbf{não altera a fonte original} ---
\texttt{map} devolve valores novos, a lista original permanece intacta; e
um stream é de \textbf{uso único} --- tentar reaproveitar o mesmo objeto
de stream depois de consumido dispara \texttt{IllegalStateException}.

Trio funcional, resumido: \texttt{filter} escolhe \textbf{quais}
elementos passam adiante (não transforma nenhum); \texttt{map}
transforma \textbf{cada} elemento, um a um; \texttt{reduce} combina todos
num \textbf{único} resultado final. Uma \textbf{lambda} é apenas a forma
curta de escrever a implementação de uma interface de método único
(\texttt{Predicate}, \texttt{Function}) no lugar de uma classe anônima
--- é açúcar de sintaxe, não um novo mecanismo de tipo.
\end{conceito}

\subsection{var, record e switch como expressão}

\begin{conceito}
\textbf{\texttt{var} (Java 10)}: infere o tipo de uma variável
\textbf{local} a partir do que a inicializa, em \textbf{tempo de
compilação} --- e o tipo inferido \textbf{fica fixo} a partir daí; Java
continua estaticamente tipado, não virou dinamicamente tipado. Exige um
inicializador (\texttt{var x;} sozinho não compila) e não se aplica a
atributo de instância, parâmetro de método nem tipo de retorno.

\textbf{\texttt{record} (Java 16)}: classe de dados imutável, declarada
por um cabeçalho compacto --- \texttt{record Ponto(int x, int y) \{ \}}.
O compilador gera automaticamente: um construtor canônico (que pode ser
sobrescrito na forma compacta para adicionar validação), métodos de
acesso no formato \texttt{p.x()}, \texttt{p.y()} (\textbf{sem} prefixo
\texttt{get}), \texttt{equals}/\texttt{hashCode} comparando componente a
componente, e \texttt{toString} no formato \texttt{Ponto[x=1, y=2]}. Os
campos de um record são \textbf{finais}; o record em si é implicitamente
\texttt{final} e já estende \texttt{java.lang.Record}, portanto
\textbf{não pode estender outra classe} --- mas \textbf{pode implementar
interfaces}. Não aceita campo de instância adicional além dos declarados
no cabeçalho (só campos estáticos).

\textbf{\texttt{switch} como expressão (Java 14)}: devolve um valor
diretamente e, na forma com seta (\texttt{->}), executa \textbf{só} o
ramo correspondente --- sem \texttt{break} e sem \textit{fall-through}
entre ramos (use \texttt{yield} dentro de um bloco com chaves quando
precisar de lógica mais elaborada num ramo).
\end{conceito}

\begin{cuidado}
"\texttt{var} torna a variável de tipo dinâmico" é falso: depois de
inferido, atribuir um valor de tipo incompatível continua sendo erro de
compilação. "O \texttt{record} gera \textit{getters} no padrão
\texttt{getX()}" é falso --- o acessor chama-se apenas \texttt{x()}, sem
prefixo. \texttt{sealed} \textbf{não é sinônimo} de \texttt{final}: o
\texttt{final} proíbe qualquer herança; o \texttt{sealed} permite herança
apenas pelos tipos explicitamente autorizados.
\end{cuidado}

\section{Leitura ativa de código em Java}
\label{sec:leitura-ativa-java-teoria}

Java é a linguagem onde erros gostam de se esconder \textbf{dentro} de
uma estrutura que parece correta à primeira leitura --- um
\texttt{try-catch} que compila mas nunca captura a exceção certa, um
escopo de variável que parece visível mas não é. O protocolo de leitura
ativa (Capítulo \ref{cap:programacao-teoria},
Seção~\ref{sec:leitura-ativa-programacao-teoria}) se aplica igualmente
aqui: releia linha a linha, sem assumir.

\subsection{Checked/unchecked disfarçadas em try-catch}

\begin{conceito}
\begin{itemize}
\item \textbf{Ordem dos catch}: um \texttt{catch (Exception e)}
\textbf{antes} de um \texttt{catch (IOException e)} \textbf{não compila}
--- o \texttt{catch} mais genérico precisa vir \textbf{por último},
porque um catch mais específico depois de um genérico nunca seria
alcançado.
\item \textbf{Catch que não cobre o tipo lançado}: um método declara
\texttt{throws SQLException}, mas o bloco que o chama só tem
\texttt{catch (RuntimeException e)} --- isso \textbf{não compila}, porque
\texttt{SQLException} é \textbf{checked} e exige tratamento explícito ou
propagação com \texttt{throws} no método chamador.
\item \textbf{\texttt{finally} que engole a exceção}: um \texttt{return}
dentro do bloco \texttt{finally} \textbf{sobrescreve} qualquer exceção
lançada no \texttt{try}/\texttt{catch} anterior --- a exceção original se
perde silenciosamente, sem rastro.
\item \textbf{Multi-catch}: \texttt{catch (IOException | SQLException e)}
exige que as duas exceções listadas \textbf{não tenham relação de
herança} entre si (uma não pode ser subtipo da outra) --- senão não
compila, por redundância.
\end{itemize}
\end{conceito}

\begin{cuidado}
Um bloco com a ordem dos \texttt{catch} invertida (genérico antes do
específico) \textbf{não compila e não funciona normalmente} --- é erro
de compilação, não um comportamento "que funciona mas é estranho". Uma
exceção \textbf{checked} lançada dentro de um método sem
\texttt{throws} e sem \texttt{try-catch} \textbf{não compila} --- não é
"ignorada silenciosamente em runtime". Erro de compilação e exceção em
tempo de execução são fenômenos diferentes, e a prova testa se essa
distinção está clara.
\end{cuidado}

\subsection{Escopo de variável}

\begin{conceito}
\begin{itemize}
\item \textbf{\textit{Shadowing} (sombreamento)}: uma variável local com o
\textbf{mesmo nome} de um atributo da classe \textbf{esconde} o atributo
dentro daquele método --- \texttt{this.atributo} continua acessível
explicitamente, mas \texttt{atributo} sozinho passa a se referir à
variável local.
\item \textbf{Variável efetivamente final em lambda/classe anônima}: uma
variável local usada dentro de uma lambda ou classe anônima precisa ser
\texttt{final} ou \textbf{efetivamente final} (nunca reatribuída depois
de inicializada) --- código que reatribui essa variável depois de criar a
lambda \textbf{não compila}.
\item \textbf{Escopo de bloco em \texttt{for}}: uma variável declarada
\textbf{dentro} do cabeçalho do \texttt{for} não existe fora dele ---
usá-la depois do laço termina é erro de compilação, não um "valor
residual" da última iteração.
\end{itemize}
\end{conceito}

\begin{cuidado}
Usar uma variável declarada dentro de um bloco (\texttt{if},
\texttt{for}, \texttt{try}) fora desse bloco é erro de compilação por
variável fora de escopo --- diferente de JavaScript, onde \texttt{var}
pode vazar para o escopo da função inteira; em Java o escopo é sempre de
\textbf{bloco}.
\end{cuidado}

\subsection{Armadilhas de threads virtuais: \textit{pinning}}

\begin{conceito}
Uma thread virtual roda sobre uma \textit{carrier thread} de plataforma.
Se o código dentro de um bloco \texttt{synchronized} faz uma chamada
\textbf{bloqueante} (I/O, banco de dados), a thread virtual pode ficar
\textbf{presa (\textit{pinned})} à sua \textit{carrier} em vez de
liberá-la para outra thread virtual usar --- isso anula justamente a
vantagem de escala das threads virtuais, porque agora uma única operação
bloqueante ocupa uma \textit{carrier} inteira, impedindo outras threads
virtuais de progredir nela. O cenário de risco típico:
\texttt{synchronized} combinado com chamada bloqueante \textbf{dentro} do
bloco sincronizado, sob alta concorrência.
\end{conceito}

\begin{cuidado}
Threads virtuais \textbf{não eliminam} toda preocupação com bloqueio ou
concorrência --- o cenário de \texttt{synchronized} + I/O bloqueante
ainda pode prender a thread virtual à sua \textit{carrier}. Elas reduzem
o \textbf{custo de criar} threads (permitindo milhares delas), não
eliminam o cuidado com seções críticas de código.
\end{cuidado}
